% =============================================================================
%  Parameter selection and error analysis.
%  SHARED: \input by both paper.tex (as Section 5) and technical-report.tex,
%  so the two documents cannot report different numbers.
% =============================================================================
\section{Parameter Selection and Error Analysis}
\label{sec:params}

The FLPSI layer is parameterised by the code length $d$, the number of sub-samples $T$, the
sub-sample width $w$, and the acceptance threshold $\theta$. This section derives the error
rates those parameters induce, and shows that the published choice
$(d,T,w,\theta) = (128, 64, 14, 2)$ is wrong by three orders of magnitude on real codes --- not
because the arithmetic is wrong, but because the distribution it assumes is.

\subsection{The accept model}

A sub-sample selects $w$ of the $d$ positions uniformly at random and collides when the two
codes agree on all of them. Writing $H = \|\xt \oplus \yt_i\|_1$ for the Hamming distance, the
per-sub-sample collision probability and the acceptance probability are
\begin{equation}
\label{eq:accept}
q(H) \;=\; \frac{\binom{d-H}{w}}{\binom{d}{w}},
\qquad
\Pr[\text{accept} \mid H] \;=\; \Pr\bigl[\mathrm{Binomial}(T, q(H)) \ge \theta\bigr].
\end{equation}
Acceptance depends on the pair \emph{only through $H$} --- not on which positions differ ---
because the mask is a uniformly random set. Consequently the false-accept and false-reject
rates are exact functionals of the two Hamming histograms,
\begin{equation}
\label{eq:rates}
\mathrm{FAR} = \sum_{H} P_{\mathrm{imp}}(H)\,\Pr[\text{accept}\mid H],
\qquad
\mathrm{FRR} = 1 - \sum_{H} P_{\mathrm{gen}}(H)\,\Pr[\text{accept}\mid H],
\end{equation}
with no Monte-Carlo error. Every number below is computed from~\eqref{eq:rates} over 1,200
held-out identities: 1,200 genuine and 2,877,600 impostor pairs.

\paragraph{Validating the model before using it.}
Equation~\eqref{eq:accept} assumes the $d$ positions are exchangeable, and Super-Bit's advantage
comes precisely from making bits \emph{dependent}, so this needs testing. Measuring $\hat q(H)$
on real code pairs (200 sub-sample draws per pair, $w=14$) gives $\hat q(H)/q(H)$ with
\textbf{median $1.002$, range $[0.919, 1.099]$} over $H \in [5,40]$. Bit dependence is real ---
mean inter-bit $|\rho| = 0.254$, and the correlation matrix has participation ratio $9.84$, so
the 128 bits behave like about ten independent decisions --- but it manifests in the distribution
\emph{of} $H$, not in $\Pr[\text{accept}\mid H]$. That isolates the defect to one place and
licenses~\eqref{eq:rates}.

\subsection{The codes are balanced but not uniform}

The parameters were selected under a model in which impostor codes are uniform, giving
$H \sim \mathrm{Binomial}(128, \tfrac12)$ with $\sigma = 5.66$. Table~\ref{tab:dist} compares
that model against measurement.

\begin{table}[htbp]
\centering
\caption{Impostor Hamming distribution: assumed versus measured over 2,877,600 held-out pairs.
The mean is exactly as predicted; nothing else is.}
\label{tab:dist}
\small
\begin{tabular}{@{}lrr@{}}
\toprule
& uniform model & measured \\
\midrule
mean $H$              & 64.0 & \textbf{64.0} \\
standard deviation    & 5.66 & \textbf{20.22} \\
minimum observed $H$  & --- (never below $\approx$40) & \textbf{0} \\
\midrule
genuine mean $H$ & --- & 8.3 (sd 5.67, max 34) \\
\bottomrule
\end{tabular}
\end{table}

The mean is \emph{exactly} right, because the quantiser's median-thresholding step balances
every bit (measured balance $\in [0.466, 0.540]$). \textbf{No first-moment check can detect this
failure}: a sanity test verifying that impostor codes ``look like coin flips on average'' passes
cleanly. The discrepancy is entirely in the second moment and the tail, where $\sigma$ is
$3.6\times$ larger than assumed.

The cause is structural rather than a defect of this particular bridge. A 128-bit Super-Bit code
is 128 projections of a 16-dimensional embedding, so it carries roughly 16 real degrees of
freedom. The bridge is \emph{faithful} --- $\mathbb{E}[H]/d = \arccos(\text{cos-sim})/\pi$ holds
--- and that faithfulness is the problem: it transmits the embedding's angular distribution
intact, including the fact that some fingers genuinely resemble one another. We expect any LSH
bridge over a learned low-dimensional embedding to behave the same way.

\begin{figure}[htbp]
\centering
\includegraphics[width=0.82\linewidth]{hamming-histograms.pdf}
\caption{Measured genuine and impostor Hamming distributions against the uniform model
(dashed). Left: linear scale, where the impostor mean sits exactly at $d/2$. Right: log scale,
where the left tail --- the part that sets the false-accept rate --- is several orders of
magnitude heavier than assumed.}
\label{fig:hist}
\end{figure}

\subsection{The consequence for the published operating point}

At $(w,\theta) = (14,2)$ with error-free sub-samples, evaluating~\eqref{eq:rates} against the
uniform model reproduces the published per-record false-accept rate of $2.954\times10^{-5}$
exactly, and $q(8) = \binom{120}{14}/\binom{128}{14} = 0.38493266$ matches the published
constant to all quoted digits. The arithmetic is sound. Against the \emph{measured}
distribution the same expressions give
\begin{center}\small
\begin{tabular}{@{}lrrr@{}}
\toprule
& uniform model & measured & ratio \\
\midrule
per-record FAR & $2.954\times10^{-5}$ & $\mathbf{4.578\times10^{-2}}$ & $\mathbf{1{,}550\times}$ \\
FRR & --- & $9.61\times10^{-3}$ & --- \\
\bottomrule
\end{tabular}
\end{center}

A per-record rate is not a security boundary. A $1{:}N$ identification query is tested against
every enrolled record, so what an attacker faces is
$\mathrm{FAR}_{\text{query}} = 1-(1-\mathrm{FAR})^{N}$ (Figure~\ref{fig:farn}):

\begin{center}\small
\begin{tabular}{@{}lrrrrrr@{}}
\toprule
$N$ & 1 & 10 & 100 & 1{,}000 & 5{,}000 & $10^{4}$ \\
\midrule
$\mathrm{FAR}_{\text{query}}$, published point & 0.046 & 0.374 & \textbf{0.991} & 1.000 & 1.000 & 1.000 \\
$\mathrm{FAR}_{\text{query}}$, collision floor & --- & --- & 0.0003 & 0.0028 & \textbf{0.0138} & 0.0274 \\
\bottomrule
\end{tabular}
\end{center}

FRR is $0.96\%$ throughout. \textbf{The published parameters describe a verification $(1{:}1)$
system, not an identification system}: at $N=1$ a $4.6\%$ false-accept rate is poor but coherent,
and by $N=100$ the query-level rate already exceeds $50\%$. The second row is the irreducible
floor derived in \S\ref{sec:floor}, which no parameter choice can go below.

\begin{figure}[htbp]
\centering
\includegraphics[width=0.52\linewidth]{far-vs-dbsize.pdf}
\caption{Query-level false-accept rate against database size. The published single-finger point
crosses $50\%$ at $N \approx 100$. Multi-finger fusion moves the curve by orders of magnitude;
retuning does not.}
\label{fig:farn}
\end{figure}

\subsection{Retuning does not repair it}

We swept 6,848 combinations of $w \in \{8,\dots,64\}$, $t \in \{0,\dots,8\}$ and
$\theta \in \{1,\dots,64\}$, keeping only points whose FAR estimate is supported by the
histogram (less than half the FAR mass drawn from bins holding fewer than ten observed pairs).
Taking the lowest-FRR point that meets each query-level target at $N = 5{,}000$:

only one target is reachable at all. A query-level rate of $10^{-1}$ is met by
$(w,t,\theta) = (56,1,40)$ --- at an FRR of \textbf{84.8\%} --- and $10^{-2}$, $10^{-3}$ and
$10^{-4}$ are unreachable by any point in the sweep with usable histogram support. The only
setting that reaches even a $10\%$ query FAR rejects $85\%$ of legitimate users. The
suggestion that $\theta$ should grow as $\Theta(\log N)$ to hold the query-level rate constant
is directionally right and quantitatively insufficient, for the reason we give next.

\subsection{An irreducible collision floor}
\label{sec:floor}

\textbf{8 of the 2,877,600 impostor pairs have Hamming distance exactly zero} --- distinct
fingers producing \emph{identical} 128-bit codes. The collision rate is
\begin{equation}
p_0 = 2.78\times10^{-6}, \qquad \text{exact Poisson 95\% CI } [1.20\times10^{-6},\, 5.48\times10^{-6}].
\end{equation}
Since $q(0) = 1$, every sub-sample collides and $\Pr[\text{accept}\mid H{=}0] = 1$ for
\emph{every} $(w, t, \theta)$. Therefore $p_0$ is a \textbf{hard lower bound on the per-record
false-accept rate of any FLPSI parameterisation over these codes}, and it propagates to the query
level as tabulated above. This is precisely why the ladder above becomes unreachable below $10^{-2}$: at $N = 5{,}000$ the
floor alone is $1.4\%$, independent of $\theta$. The near tail is dense as well --- $1.03\times
10^{-3}$ of impostor pairs lie at $H \le 10$, and $2.80\times10^{-2}$ at $H \le 25$, which is the
99th percentile of the \emph{genuine} distribution. The two distributions genuinely overlap; the
protocol is being asked to separate things that are not separable.

\textbf{The bound is a property of the code, not of the protocol.} It moves only by producing
better codes, and it is cheap for any future system to report: count $H = 0$ pairs on a held-out
set. We recommend it as a standard diagnostic for fuzzy-PSI-over-learned-codes constructions.

\subsection{Multi-finger fusion is the lever that works}

Each finger is a separate enrolment, so $k$-of-$K$ fusion multiplies evidence instead of trading
FAR against FRR along one curve --- and it is the only mechanism here that can beat the
single-finger floor, since $k$ fingers must collide simultaneously. We measure it over 201
held-out subjects contributing at least three fingers each, computing a Poisson-binomial over
the real per-finger accept probabilities, so any correlation between one person's own fingers is
carried by the data rather than assumed away.

\begin{table}[htbp]
\centering
\caption{Measured $k$-of-$K$ fusion at $N = 5{,}000$, over 201 held-out subjects.}
\label{tab:fusion}
\small
\begin{tabular}{@{}llrrr@{}}
\toprule
parameters & rule & per-record FAR & $\mathrm{FAR}_{\text{query}}$ & FRR \\
\midrule
$w{=}14$, $\theta{=}2$ & 1-of-3 & $1.31\times10^{-1}$ & 1.000 & 0.000 \\
$w{=}14$, $\theta{=}2$ & 2-of-3 & $8.73\times10^{-3}$ & 1.000 & 0.001 \\
$w{=}14$, $\theta{=}2$ & 3-of-3 & $3.19\times10^{-4}$ & 0.797 & 0.026 \\
$w{=}24$, $\theta{=}2$ & 1-of-3 & $1.95\times10^{-2}$ & 1.000 & 0.002 \\
$w{=}24$, $\theta{=}2$ & 2-of-3 & $2.45\times10^{-4}$ & 0.707 & 0.033 \\
\textbf{$w{=}24$, $\theta{=}2$} & \textbf{3-of-3} & $\mathbf{1.30\times10^{-6}}$ &
  \textbf{0.0065} & \textbf{0.281} \\
\bottomrule
\end{tabular}
\end{table}

Fusion buys four orders of magnitude on the per-record rate, far more than any single-finger
retuning achieves. It is still not sufficient on its own: the best measured configuration costs
a $28\%$ false-reject rate. The honest summary is that \textbf{at $N = 5{,}000$ with a
16-dimensional embedding, no configuration of this protocol is simultaneously secure and
usable}.

\paragraph{What this means for parameter selection.} Report FLPSI biometric parameters as a
verification operating point unless the query-level rate at the target $N$ is stated; fit them to
the \emph{measured} code distribution, since the uniform assumption errs by $1{,}550\times$ here
and fails silently; and publish the collision floor of \S\ref{sec:floor}, which bounds every
parameter choice and costs one pass over a held-out set. The remedy is more evidence, not a
different threshold, and the binding constraint is the embedding dimension --- sixteen dimensions
cannot separate 5,000 identities at cryptographic error rates.
